\section{Methods}\label{sec:methods}

\subsection{Jets, classes and benchmarks}
A jet is the detector signature of an energetic quark or gluon, or of a heavy particle whose decay products are
collimated by its momentum.
For a machine-learning model it is an unordered set of a few tens of reconstructed particles, each described here
by \FillNFeatures{} numbers: its momentum relative to the jet, its electric charge and type and, for charged
particles, the displacement of its track from the collision point, which reveals the decays of the long-lived
hadrons that contain bottom or charm quarks.
The classes of JetClass-II name what the jet came from: either a resonance, a short-lived heavy particle labeled
by the quarks, gluons and leptons into which it decayed, each of which appears as one prong of the jet, or the
ordinary production of quarks and gluons by the strong interaction
(QCD jets), which is the background in most analyses.
The two binary benchmarks are standard in the field: top tagging asks whether a jet contains the decay of a top
quark, and quark--gluon discrimination whether it was initiated by a quark or by a gluon.

\subsection{Data and model}
We pretrain on JetClass-II~\cite{sophon,jetclasstwo}, the dataset of the Signature-Oriented Pre-training for
Heavy-resonance ObservatioN (Sophon) model, whose \Kfull{} classes are \LitJetClassTwoNResonantClasses{}
resonant classes and \LitJetClassTwoNQcdClasses{} QCD classes.
A resonant jet contains the decay of a hypothetical particle $X$, either directly into two quarks, gluons or
leptons or through two intermediate particles into three or four of them, and a QCD class is defined by the
number of bottom, charm and strange quarks in the jet.
The jets are large, with a radius of \LitJetClassTwoJetRadius, and come from a simulated detector with pileup,
the additional collisions that overlap each event; after Sophon's selection there are \DesignNTrainSelected{}
training jets (appendix~\ref{app:training}).
In pretraining the classes keep the frequencies they have in the dataset, which we call the natural mix, and
jets are drawn with Sophon's weights, which give each of its sampling groups of classes the same distribution of
$p_{\mathrm T}$ and mass.

The network is Sophon's~\cite{sophon}, a Particle Transformer~\cite{part} with \DesignParticleBlocks{} attention
blocks over the particles, \DesignClassBlocks{} blocks that pool them into one class token of
\DesignEmbedDim{} numbers, and an output layer with one hidden layer of \DesignHeadHidden{} units.
The class token is what we call the frozen features, the input of the output layer, which every probe reads.


\subsection{Label sets}
A label set is a map from the \Kfull{} classes onto fewer classes, and a model pretrained with it learns to
predict the merged class of each jet.
The four label sets are nested, so that each is obtained from the one before by merging classes that differ in
a named way (table~\ref{tab:sets} and figure~\ref{fig:design}; the full map is table~\ref{tab:classmap}).
The \Kres-class set merges the QCD classes into one and keeps every resonant class.
\input{tables/table1}%
\vfig[13cm]{fig1_design}{Design of the study. (a) The ten classes of one group of the \Kcoarse-class set, a resonance $X$
decaying to two quarks or gluons, drawn as schematic jets, with the class each falls in under the four label sets;
below, all \Kfull{} classes as cells, where a bar spanning several cells is one class of that label set. (b) One model
is pretrained for each label set on the same jets and is then used with frozen features or fine-tuned.}{fig:design}
The \Kmid-class set then labels a decay only by whether it contains a bottom quark, else a charm quark, else
neither, and no longer separates strange from up and down quarks, gluons from these light quarks, or electrons
from muons.
The \Kcoarse-class set removes the remaining flavor and keeps what every set keeps, the number of decay products
and whether each is a quark or gluon, an electron or muon, a tau lepton or a neutrino.
Because the sampling groups do not depend on the label set, models pretrained with the same random seed see the
same jets in the same order, and they start from the same weights in every layer except the last, whose width
is the number of classes.

A label set also limits what the output layer can provide.
For a classifier at the minimum of the cross-entropy, the score $g_c$ of a merged class $c$, its softmax output,
is the sum of the scores $g'_i$ that the \Kfull-class model gives the classes it contains,
\begin{equation}
  g_c(x) = \sum_{i\,\in\, c} g'_i(x) ,
  \label{eq:sum}
\end{equation}
so a physics analysis can form its discriminant, the quantity on which it selects jets, by summing scores only
if the classes it needs are unions of classes of the label set~\cite{sophon}.


\subsection{Random groupings and mass regression}
The random groupings test whether the number of classes or their content matters.
Each assigns the \LitJetClassTwoNResonantClasses{} resonant classes at random to \RandVtwoNResGroups{} groups
whose shares of the training jets match those of the \Kcoarse-class set to within \RandVtwoShareTol, and keeps
the QCD jets as one class.
The groupings are not typical random draws: among draws that satisfy the share rule we selected those that
merge the fewest pairs of classes that the \Kcoarse-class set also merges.


The mass regression adds a target that no label set contains.
The \Kres- and \Kcoarse-class models are pretrained a second time with one more output $\hat\theta$, trained
to predict $\theta = \ln(m_{\mathrm{gen}}/m_{\mathrm{SD}})$, where $m_{\mathrm{SD}}$ is the reconstructed jet mass, computed after soft, wide-angle particles are removed
(the soft-drop mass~\cite{softdrop}), and $m_{\mathrm{gen}}$ the same mass before detector effects, with the
loss
\begin{equation}
  L = L_{\mathrm{cls}} + \lambda \log\cosh\bigl(\hat\theta(x) - \theta\bigr) ,
  \label{eq:massloss}
\end{equation}
where $L_{\mathrm{cls}}$ is the cross-entropy and $\lambda = \MlstMassLambda$, the weight used by the CMS
taggers that introduced this term~\cite{glopart}.

\subsection{Evaluation}
Linear probes measure what the frozen features hold.
A probe is an L2-regularized multinomial logistic regression on the standardized features.
On JetClass-II the probes are trained and tested on approximately class-balanced sets of jets that no
pretraining run read, so that chance is $1/\Kfull$, or \VChance\%.
Binary probes on pairs of JetClass-II classes ask whether a label set preserves one particular distinction, such
as $X\to b\bar b$ against $X\to c\bar c$.
Two references are probed like the pretrained models: an untrained network, with the statistics of its input
normalization computed on training jets, and the checkpoint released with Sophon~\cite{sophon}.

Fine-tuning measures what survives when every weight adapts.
The pretrained output layer is replaced by a new one of the same shape with random weights, and the same network
trained from random weights on the downstream task alone, with \VScratchRunsWord{} seeds, is the reference.
On JetClass-II and JetClass we fine-tune on \FtSizeESix{} jets, with one learning rate for all pretrained
models, chosen on validation for each task and training size, and on top tagging and quark--gluon discrimination
on the full training sets with the published recipe of the Particle Transformer~\cite{part}.
The fine-tuned models are scored on the test jets of the probes.

The tasks differ in how far they lie from pretraining.
JetClass~\cite{part} has ten classes of jets from Higgs bosons, $W$ and $Z$ bosons, top quarks and QCD; it
was simulated without pileup and in a narrower range of transverse momentum.
The top-tagging~\cite{toptag} and quark--gluon~\cite{efn,qgpythia} data come from other simulations and lack
inputs that the network was pretrained with: the first provides only the momenta of the particles,
\FillTopNPresent{} of the \FillNFeatures{} inputs, and the second lacks the track displacements
(appendix~\ref{app:inputs}).
In fine-tuning the absent inputs are set to zero, and for the probes, which cannot adapt, they are filled with
values chosen on validation.
A frozen result on these two benchmarks is thus a linear readout of \DesignEmbedDim{} numbers computed from
incomplete inputs.

We report the accuracy, the area under the receiver operating characteristic curve (AUC), averaged over
one-versus-rest problems when a task has more than two classes, and the background rejection
$\mathrm{Rej}_{30\%}$, the inverse of the false-positive rate at a true-positive rate of $30\%$~\cite{part};
$\mathrm{Rej}_{50\%}$ is the same at $50\%$.

The mass resolution measures how well the jet mass can be recovered.
For an estimate $m_{\mathrm{est}}$ and the response $r = m_{\mathrm{est}}/m_{\mathrm{gen}}$ we use
\begin{equation}
  \sigma_r = \frac{q_{84.1} - q_{15.9}}{2\,q_{50}} ,
  \label{eq:resolution}
\end{equation}
with $q_p$ the $p$th percentile of $r$.
We compare three estimates on the JetClass-II test jets of every class for which $m_{\mathrm{gen}}$ exists: the
reconstructed soft-drop mass, the regression output of the models trained with equation~(\ref{eq:massloss}), and
a multilayer-perceptron (MLP) probe that predicts $\theta$ from the frozen features of any model, the last two
as $m_{\mathrm{est}} = m_{\mathrm{SD}}\,e^{\hat\theta}$.

The resonance benchmark repeats the test with which Sophon was introduced~\cite{sophon}: a particle $X$ of mass
\XbxMx~GeV decaying to $b\bar b$ or to $bs$ is to be separated from QCD jets with a mass between \XbxWindowLo{}
and \XbxWindowHi~GeV, and we quote the signal efficiency at a background efficiency of \VXbxEpsB\%.
Sophon selects $X\to b\bar b$ with a discriminant built from its output scores, and $X\to bs$, which is not among
its classes, by what it calls transfer learning: a network with one hidden layer, started from that of the
pretrained output layer, is trained on the frozen features (appendix~\ref{app:training}).
We use the transfer learning for both decays and every label set, since $X\to b\bar b$ is a class of its own
only in the two finest, and compare it with the same network trained from random weights.

\subsection{Uncertainties}
Tables give the mean and standard deviation over independent pretraining runs, each followed by one probe or
one fine-tuning: \VRunsWord{} runs for each nested label set, \VMassRunsWord{} for each model with mass
regression, and one for each of the \VRandWord{} random groupings, whose spread therefore mixes the difference
between groupings with that between runs.
A difference $\Delta$ between the means of two models is called significant when $|\Delta|$ exceeds three times
$\sigma$, where $\sigma^2$ adds the squared standard deviations over runs of both models, not standard errors,
and the variance of $\Delta$ over bootstrap resamples of the test jets, drawn once for both models.
With so few runs this is a convention, not a Gaussian probability, and we make no correction for the number of
comparisons.
The tested differences, with each term of $\sigma$, are in the data release.
